\documentclass[11pt]{article}
\usepackage[margin=1in]{geometry}
\usepackage{booktabs}
\usepackage{enumitem}
\usepackage{hyperref}
\usepackage{longtable}
\usepackage{array}
\setlength{\parindent}{0pt}
\setlength{\parskip}{0.6em}

\title{Prompt-Profile Combined Audit:\\Natural Regression, Balanced Binary, and Metadata Controls}
\author{}
\date{2026-04-05 02:10 UTC}

\begin{document}
\maketitle

\section*{Executive Summary}
\begin{itemize}[leftmargin=1.5em]
\item This report bundles the current prompt-profile surface into one note:
  \begin{itemize}[leftmargin=1.5em]
  \item canonical regression on \texttt{mean\_relative\_length}
  \item current binary recommendation on \texttt{majority\_s\_0.5}
  \item a new metadata-baseline audit over the saved prompt archives
  \item an Athena code audit over the prompt-profile build / label / train / summarize path
  \end{itemize}
\item Regression is now pinned cleanly.
  \begin{itemize}[leftmargin=1.5em]
  \item Slurm \texttt{2215} reran the canonical natural regression object from the current PR head.
  \item It reproduced the original locked \texttt{2043} regression ledger exactly.
  \item Max absolute difference was \texttt{0.0} across the prompt-length baselines plus \texttt{ensemble} and \texttt{last\_layer} at both frozen checkpoints.
  \end{itemize}
\item Binary is still the cleaner deployment-facing head.
  \begin{itemize}[leftmargin=1.5em]
  \item The current best single global activation surface remains \texttt{ensemble h256 d1}.
  \item Cross-dataset mean test \texttt{PR-AUC}: prompt length \texttt{0.350}, \texttt{ensemble h256 d1} \texttt{0.518}, \texttt{last\_layer h256 d2} \texttt{0.430}.
  \end{itemize}
\item The metadata story is narrower than the older April PDFs.
  \begin{itemize}[leftmargin=1.5em]
  \item The current reported ``metadata baseline'' is really a train-fit 1D prompt-length baseline, because \texttt{effective\_max\_tokens=30000} is fixed.
  \item Cheap prompt-shape features often match or beat prompt length, and in a few cases rival or beat the activation probes too.
  \item So the current results establish lift over prompt length on some surfaces, not yet lift over strong prompt-only controls in general.
  \end{itemize}
\end{itemize}

\section*{Current Objects}
\subsection*{Regression}
\begin{itemize}[leftmargin=1.5em]
\item target: \texttt{mean\_relative\_length}
\item split: natural prompt-disjoint train/test
\item sampler: natural
\item views: \texttt{ensemble}, \texttt{last\_layer}
\item canonical artifact:
  \begin{itemize}[leftmargin=1.5em]
  \item \texttt{docs/weeks/2026-W14/prompt-profile-natural-regression-rerun-2026-04-05.md}
  \item \texttt{outputs/prompt\_profile\_natural\_regression\_rerun\_20260405/prompt\_profile\_natural\_regression\_rerun\_20260405.pdf}
  \end{itemize}
\end{itemize}

\subsection*{Binary}
\begin{itemize}[leftmargin=1.5em]
\item target: \texttt{majority\_s\_0.5}
\item split: train balanced by downsampling, test natural
\item current recommended activation surface: \texttt{ensemble h256 d1}
\item capacity-control artifact:
  \begin{itemize}[leftmargin=1.5em]
  \item \texttt{docs/weeks/2026-W14/prompt-profile-binary-capacity-controls-2026-04-04.md}
  \item \texttt{outputs/prompt\_profile\_binary\_capacity\_controls\_20260404/prompt\_profile\_binary\_capacity\_controls\_20260404.pdf}
  \end{itemize}
\end{itemize}

\section*{Regression Read}
Primary metric here is screening, not calibration.

\begin{center}
\small
\begin{tabular}{@{}lrrrr@{}}
\toprule
Dataset & Prompt length \texttt{top\_20p\_capture} & Ensemble \texttt{top\_20p\_capture} & Prompt length \texttt{RMSE} & Ensemble \texttt{RMSE} \\
\midrule
\texttt{GPQA} & 0.168 & 0.247 $\pm$ 0.018 & 0.165 & 0.160 $\pm$ 0.002 \\
\texttt{AIME} & 0.306 & 0.305 $\pm$ 0.017 & 0.176 & 0.183 $\pm$ 0.006 \\
\texttt{MATH-500} & 0.290 & 0.262 $\pm$ 0.058 & 0.154 & 0.162 $\pm$ 0.002 \\
\texttt{MMLU-Pro} & 0.292 & 0.355 $\pm$ 0.025 & 0.126 & 0.133 $\pm$ 0.000 \\
\texttt{LiveCodeBench} & 0.248 & 0.340 $\pm$ 0.005 & 0.279 & 0.248 $\pm$ 0.018 \\
\bottomrule
\end{tabular}
\end{center}

Cross-dataset means:
\begin{itemize}[leftmargin=1.5em]
\item \texttt{top\_20p\_capture}: prompt length \texttt{0.261}, ensemble \texttt{0.302}
\item \texttt{RMSE}: prompt length \texttt{0.180}, ensemble \texttt{0.177}
\end{itemize}

Interpretation:
\begin{itemize}[leftmargin=1.5em]
\item On the operational screening read, ensemble beats prompt length on \texttt{GPQA}, \texttt{MMLU-Pro}, and \texttt{LiveCodeBench}.
\item On calibration, prompt length still wins on \texttt{AIME}, \texttt{MATH-500}, and \texttt{MMLU-Pro}.
\item So this lane is still useful as a prompt-conditioned long-completion risk score.
\item It is still not a clean calibrated regressor, and it is not yet evidence of a loop-specific hidden-state signal by itself.
\end{itemize}

\section*{Binary Read}
\begin{center}
\small
\begin{tabular}{@{}lrrrr@{}}
\toprule
Dataset & Test prevalence & Prompt length \texttt{PR-AUC} & Ensemble \texttt{h256 d1} & Last-layer \texttt{h256 d2} \\
\midrule
\texttt{GPQA} & 0.0500 & 0.0657 & 0.5833 & 0.3176 \\
\texttt{AIME} & 0.5833 & 0.8976 & 0.9120 & 0.8479 \\
\texttt{MATH-500} & 0.0400 & 0.1002 & 0.1663 & 0.1322 \\
\texttt{MMLU-Pro} & 0.0125 & 0.1104 & 0.2116 & 0.0959 \\
\texttt{LiveCodeBench} & 0.3375 & 0.5760 & 0.7143 & 0.7575 \\
\bottomrule
\end{tabular}
\end{center}

Cross-dataset mean test \texttt{PR-AUC}:
\begin{itemize}[leftmargin=1.5em]
\item prompt length \texttt{0.350}
\item \texttt{ensemble h256 d1} \texttt{0.518}
\item \texttt{last\_layer h256 d2} \texttt{0.430}
\end{itemize}

Interpretation:
\begin{itemize}[leftmargin=1.5em]
\item \texttt{majority\_s\_0.5} is still the cleaner finished head.
\item \texttt{ensemble h256 d1} remains the best single global activation surface.
\item But the prompt-length baseline is too strong on \texttt{AIME} to support a broad ``activation-only lift'' story there.
\item \texttt{LiveCodeBench} and \texttt{GPQA} remain the cleanest current binary activation wins.
\end{itemize}

\section*{Metadata Audit}
\subsection*{What the Current Reported Baseline Really Is}
\begin{itemize}[leftmargin=1.5em]
\item The main report baseline is intentionally narrow: a train-fit prompt-only scorer on \texttt{prompt\_token\_count} and \texttt{effective\_max\_tokens}.
\item On this fixed-budget surface, \texttt{effective\_max\_tokens=30000} is constant.
\item So the effective control is just prompt length.
\item That means the current report does not yet compare activations against a strong prompt-only metadata ceiling.
\end{itemize}

\subsection*{Cheap Prompt-Stat Audit}
The new audit used the saved prompt archives directly:
\begin{itemize}[leftmargin=1.5em]
\item regression source root: the reused shared archives behind
  \texttt{outputs/prompt\_profile\_natural\_regression\_rerun\_20260405/}
\item binary source root:
  \texttt{outputs/prompt\_profile\_full\_train\_locked\_pair\_20260404/}
\item saved audit bundle:
  \texttt{outputs/prompt\_profile\_metadata\_audit\_20260405/}
\end{itemize}

\begin{center}
\small
\begin{tabular}{@{}>{\raggedright\arraybackslash}p{1.6in}>{\raggedright\arraybackslash}p{2.1in}>{\raggedright\arraybackslash}p{2.1in}@{}}
\toprule
Dataset & Regression best raw feature by \texttt{top\_20p\_capture} & Binary best raw feature by \texttt{PR-AUC} \\
\midrule
\texttt{GPQA} & \texttt{newline\_count 0.242} & \texttt{newline\_count 0.167} \\
\texttt{AIME} & \texttt{char\_length 0.332} & \texttt{dollar\_count 0.937} \\
\texttt{MATH-500} & \texttt{prompt\_token\_count 0.290} & \texttt{dollar\_count 0.117} \\
\texttt{MMLU-Pro} & \texttt{dollar\_count 0.300} & \texttt{newline\_count 0.506} \\
\texttt{LiveCodeBench} & \texttt{prompt\_token\_count 0.248} & \texttt{char\_length 0.590} \\
\bottomrule
\end{tabular}
\end{center}

Important comparisons against raw prompt length:
\begin{itemize}[leftmargin=1.5em]
\item \texttt{AIME} binary:
  \begin{itemize}[leftmargin=1.5em]
  \item prompt length \texttt{PR-AUC 0.898}
  \item dollar count \texttt{PR-AUC 0.937}
  \item current recommended activation surface \texttt{ensemble h256 d1 0.912}
  \end{itemize}
\item \texttt{MMLU-Pro} binary:
  \begin{itemize}[leftmargin=1.5em]
  \item prompt length \texttt{PR-AUC 0.110}
  \item newline count \texttt{PR-AUC 0.506}
  \item current recommended activation surface \texttt{ensemble h256 d1 0.212}
  \end{itemize}
\item \texttt{GPQA} regression:
  \begin{itemize}[leftmargin=1.5em]
  \item prompt length \texttt{top\_20p\_capture 0.168}
  \item newline count \texttt{0.242}
  \item current ensemble \texttt{0.247}
  \end{itemize}
\end{itemize}

So the cheap-feature audit already shows that the current 1D prompt-length control is missing real prompt-shape structure.

\subsection*{Prompt-Length Quartiles}
The quartile picture is also informative:
\begin{itemize}[leftmargin=1.5em]
\item \texttt{AIME}:
  \begin{itemize}[leftmargin=1.5em]
  \item prompt-length quartiles rise almost monotonically in both \texttt{mean\_relative\_length} and \texttt{majority\_s\_0.5}
  \item the longest quartile reaches mean relative length \texttt{0.765} and binary positive rate \texttt{1.000}
  \end{itemize}
\item \texttt{MMLU-Pro}:
  \begin{itemize}[leftmargin=1.5em]
  \item the longest quartile moves from mean relative length \texttt{0.075} to \texttt{0.183}
  \item binary positive rate moves from \texttt{0.000} to \texttt{0.050}
  \end{itemize}
\item \texttt{LiveCodeBench}:
  \begin{itemize}[leftmargin=1.5em]
  \item the middle-to-long quartiles are the riskiest
  \item binary positive rate rises from \texttt{0.100} in the shortest quartile to \texttt{0.538} in the third quartile
  \end{itemize}
\item \texttt{GPQA}:
  \begin{itemize}[leftmargin=1.5em]
  \item the effect is non-monotone
  \item the third quartile is riskiest, and the very longest prompts get safer again
  \item this argues for prompt-family structure, not a simple linear length effect
  \end{itemize}
\end{itemize}

\section*{Athena Audit}
Athena reviewed the current prompt-profile code path and the cheap prompt-stat audit in deep mode.

Files reviewed included:
\begin{itemize}[leftmargin=1.5em]
\item \texttt{docs/weeks/2026-W14/prompt-profile-natural-regression-rerun-2026-04-05.md}
\item \texttt{docs/weeks/2026-W14/prompt-profile-full-train-results-2026-04-04.md}
\item \texttt{scripts/run\_prompt\_profile\_full\_train.py}
\item \texttt{scripts/summarize\_prompt\_profile\_full\_train.py}
\item \texttt{scripts/train\_probe.py}
\item \texttt{scripts/train\_metadata\_residual\_probe.py}
\item \texttt{src/loop\_probe/labeling.py}
\item \texttt{scripts/audit\_prompt\_metadata\_correlations.py}
\item \texttt{outputs/prompt\_profile\_metadata\_audit\_20260405/metadata\_correlation\_summary.json}
\end{itemize}

Main Athena conclusions:
\begin{itemize}[leftmargin=1.5em]
\item The strongest explanation is conceptual, not a hidden implementation bug.
\item The present targets are close to fixed-budget long-completion risk:
  \begin{itemize}[leftmargin=1.5em]
  \item \texttt{mean\_relative\_length} is output length scaled by a constant budget.
  \item \texttt{majority\_s\_0.5} is effectively ``does this prompt usually cross about half the budget?''
  \end{itemize}
\item So prompt-only structure can be strong even before activations enter.
\item The missing control is a stronger prompt-shape baseline, not another prompt-length row:
  \begin{itemize}[leftmargin=1.5em]
  \item token length, log token length, char length
  \item newline, digit, dollar counts
  \item prompt-template / subject / family identity where available
  \item nonlinearities or interactions, because \texttt{GPQA} is already visibly non-monotone
  \end{itemize}
\item The current results still support narrower claims:
  \begin{itemize}[leftmargin=1.5em]
  \item ensemble remains better than \texttt{last\_layer} on most current surfaces
  \item \texttt{ensemble h256 d1} is still the right current binary default
  \item \texttt{LiveCodeBench} remains the clearest regression and binary activation win
  \end{itemize}
\item What they do not yet support is a general claim of hidden-state lift over strong prompt-only controls, or a clean loop-specific early-warning claim.
\end{itemize}

\section*{Recommended Next Checks}
Priority order from the combined local audit plus Athena:
\begin{enumerate}[leftmargin=1.5em]
\item Replace the current 1D prompt-length control with a strong train-fit prompt-shape baseline on the same frozen splits.
\item Evaluate activation probes as residuals over that metadata model, not only as stand-alone predictors.
\item Re-evaluate inside matched strata of prompt shape:
  prompt-length quantile, newline / dollar / digit bins, and subject or template family where available.
\item Use family-held-out splits where possible, not only prompt-disjoint random splits.
\item Decompose ``long completion risk'' from ``loop-specific risk'' more explicitly before making stronger mechanistic claims.
\end{enumerate}

\section*{Current Honest Claim}
The combined current claim should be:
\begin{itemize}[leftmargin=1.5em]
\item Regression:
  \begin{itemize}[leftmargin=1.5em]
  \item useful as a screening score for prompt-conditioned long-completion risk
  \item still mixed against even the narrow prompt-length baseline
  \end{itemize}
\item Binary:
  \begin{itemize}[leftmargin=1.5em]
  \item \texttt{majority\_s\_0.5} is the cleaner finished head
  \item \texttt{ensemble h256 d1} is the best current global activation surface
  \end{itemize}
\item Metadata:
  \begin{itemize}[leftmargin=1.5em]
  \item the current report establishes lift over 1D prompt length on some surfaces
  \item it does not yet establish lift over strong prompt-only controls in general
  \end{itemize}
\end{itemize}

\section*{Artifacts}
\begin{itemize}[leftmargin=1.5em]
\item combined audit PDF:
  \texttt{outputs/prompt\_profile\_combined\_audit\_20260405/prompt\_profile\_combined\_audit\_20260405.pdf}
\item combined audit TeX:
  \texttt{outputs/prompt\_profile\_combined\_audit\_20260405/prompt\_profile\_combined\_audit\_20260405.tex}
\item metadata audit bundle:
  \texttt{outputs/prompt\_profile\_metadata\_audit\_20260405/}
\end{itemize}

\end{document}
